\documentclass[12pt,a4paper]{article}
\usepackage[utf8]{inputenc}
\usepackage[indonesian]{babel}
\usepackage{amsmath}
\usepackage{amsfonts}
\usepackage{amssymb}
\usepackage{geometry}

\geometry{margin=2.5cm}

\title{Menghitung Volume Benda Putar dengan Metode Kulit Tabung}
\author{Ahmad Fakhrul Bawani}
\date{}

\begin{document}

\maketitle

\section*{Soal}

Use the shell method to find the volume of the solid generated by revolving the region bound by $y = x^3$, $y = 8$, and $x = 0$ about the following lines:
\begin{center}
  \textbf{a.} The $y$-axis \quad \textbf{b.} The line $x = 5$ \quad \textbf{c.} The line $x = -2$ \\
  \textbf{d.} The $x$-axis \quad \textbf{e.} The line $y = 9$ \quad \textbf{f.} The line $y = -1$
\end{center}

\subsection*{Karakteristik Daerah Batas}
\begin{itemize}
  \item Batas dalam variabel $x$: Daerah membentang dari $x = 0$ hingga titik potong $x^3 = 8 \implies x = 2$. Tinggi elemen vertikal adalah $8 - x^3$.
  \item Batas dalam variabel $y$: Daerah membentang dari $y = 0$ hingga $y = 8$. Panjang elemen horizontal adalah $x = y^{\frac{1}{3}}$.
\end{itemize}

\subsection*{a. Diputar terhadap Sumbu-$y$ (Poros Vertikal)}
\begin{itemize}
  \item Jari-jari kulit tabung: $r = x$
  \item Tinggi kulit tabung: $h = 8 - x^3$
  \item \textbf{Setup Integral:}
    \begin{equation*}
      V = \int_{0}^{2} 2\pi \cdot x \cdot (8 - x^3) \, dx = 2\pi \int_{0}^{2} (8x - x^4) \, dx
    \end{equation*}
  \item \textbf{Hasil Perhitungan:}
    \begin{equation*}
      V = 2\pi \left[ 4x^2 - \frac{1}{5}x^5 \right]_{0}^{2} = 2\pi \left( 16 - \frac{32}{5} \right) = 2\pi \left( \frac{48}{5} \right) = \frac{96\pi}{5}
    \end{equation*}
\end{itemize}

\subsection*{b. Diputar terhadap Garis $x = 5$ (Poros Vertikal)}
\begin{itemize}
  \item Jari-jari kulit tabung: $r = 5 - x$
  \item Tinggi kulit tabung: $h = 8 - x^3$
  \item \textbf{Setup Integral:}
    \begin{equation*}
      V = \int_{0}^{2} 2\pi \cdot (5 - x) \cdot (8 - x^3) \, dx = 2\pi \int_{0}^{2} (40 - 5x^3 - 8x + x^4) \, dx
    \end{equation*}
  \item \textbf{Hasil Perhitungan:}
    \begin{align*}
      V &= 2\pi \left[ 40x - \frac{5}{4}x^4 - 4x^2 + \frac{1}{5}x^5 \right]_{0}^{2} \\
      &= 2\pi \left( 80 - 20 - 16 + \frac{32}{5} \right) = 2\pi \left( 44 + \frac{32}{5} \right) = \frac{504\pi}{5}
    \end{align*}
\end{itemize}

\subsection*{c. Diputar terhadap Garis $x = -2$ (Poros Vertikal)}
\begin{itemize}
  \item Jari-jari kulit tabung: $r = x - (-2) = x + 2$
  \item Tinggi kulit tabung: $h = 8 - x^3$
  \item \textbf{Setup Integral:}
    \begin{equation*}
      V = \int_{0}^{2} 2\pi \cdot (x + 2) \cdot (8 - x^3) \, dx = 2\pi \int_{0}^{2} (8x - x^4 + 16 - 2x^3) \, dx
    \end{equation*}
  \item \textbf{Hasil Perhitungan:}
    \begin{align*}
      V &= 2\pi \left[ 4x^2 - \frac{1}{5}x^4 - \frac{1}{2}x^4 + 16x \right]_{0}^{2} \text{ (setelah disederhanakan)} \\
      &= 2\pi \left[ 16x + 4x^2 - \frac{1}{2}x^4 - \frac{1}{5}x^5 \right]_{0}^{2} \\
      &= 2\pi \left( 32 + 16 - 8 - \frac{32}{5} \right) = 2\pi \left( 40 - \frac{32}{5} \right) = \frac{336\pi}{5}
    \end{align*}
\end{itemize}

\subsection*{d. Diputar terhadap Sumbu-$x$ (Poros Horizontal)}
\begin{itemize}
  \item Jari-jari kulit tabung: $r = y$
  \item Tinggi kulit tabung: $h = y^{\frac{1}{3}}$
  \item \textbf{Setup Integral:}
    \begin{equation*}
      V = \int_{0}^{8} 2\pi \cdot y \cdot y^{\frac{1}{3}} \, dy = 2\pi \int_{0}^{8} y^{\frac{4}{3}} \, dy
    \end{equation*}
  \item \textbf{Hasil Perhitungan:}
    \begin{equation*}
      V = 2\pi \left[ \frac{3}{7}y^{\frac{7}{3}} \right]_{0}^{8} = 2\pi \cdot \frac{3}{7} (128) = \frac{768\pi}{7}
    \end{equation*}
\end{itemize}

\subsection*{e. Diputar terhadap Garis $y = 9$ (Poros Horizontal)}
\begin{itemize}
  \item Jari-jari kulit tabung: $r = 9 - y$
  \item Tinggi kulit tabung: $h = y^{\frac{1}{3}}$
  \item \textbf{Setup Integral:}
    \begin{equation*}
      V = \int_{0}^{8} 2\pi \cdot (9 - y) \cdot y^{\frac{1}{3}} \, dy = 2\pi \int_{0}^{8} (9y^{\frac{1}{3}} - y^{\frac{4}{3}}) \, dy
    \end{equation*}
  \item \textbf{Hasil Perhitungan:}
    \begin{align*}
      V &= 2\pi \left[ 9 \cdot \frac{3}{4}y^{\frac{4}{3}} - \frac{3}{7}y^{\frac{7}{3}} \right]_{0}^{8} = 2\pi \left[ \frac{27}{4}(16) - \frac{3}{7}(128) \right] \\
      &= 2\pi \left( 108 - \frac{384}{7} \right) = 2\pi \left( \frac{756 - 384}{7} \right) = \frac{744\pi}{7}
    \end{align*}
\end{itemize}

\subsection*{f. Diputar terhadap Garis $y = -1$ (Poros Horizontal)}
\begin{itemize}
  \item Jari-jari kulit tabung: $r = y - (-1) = y + 1$
  \item Tinggi kulit tabung: $h = y^{\frac{1}{3}}$
  \item \textbf{Setup Integral:}
    \begin{equation*}
      V = \int_{0}^{8} 2\pi \cdot (y + 1) \cdot y^{\frac{1}{3}} \, dy = 2\pi \int_{0}^{8} (y^{\frac{4}{3}} + y^{\frac{1}{3}}) \, dy
    \end{equation*}
  \item \textbf{Hasil Perhitungan:}
    \begin{align*}
      V &= 2\pi \left[ \frac{3}{7}y^{\frac{7}{3}} + \frac{3}{4}y^{\frac{4}{3}} \right]_{0}^{8} = 2\pi \left[ \frac{3}{7}(128) + \frac{3}{4}(16) \right] \\
      &= 2\pi \left( \frac{384}{7} + 12 = \frac{384 + 84}{7} \right) = \frac{936\pi}{7}
    \end{align*}
\end{itemize}

\end{document}